\documentclass[lualatex,ja=standard,a4paper,12pt]{bxjsarticle}
\usepackage{amsmath, amssymb, amsthm}
\usepackage{geometry}
\usepackage{xcolor}
\usepackage{tcolorbox}
\usepackage{booktabs}
\usepackage{array}
\usepackage{graphicx}
\usepackage{tikz}
\usetikzlibrary{arrows.meta, positioning}

\geometry{margin=25mm}

\tcbuselibrary{theorems, skins, breakable}

\newtcolorbox{defbox}[1]{
  colback=blue!5, colframe=blue!60!black,
  fonttitle=\bfseries, title={#1}, arc=3pt, breakable
}
\newtcolorbox{thmbox}[1]{
  colback=green!5, colframe=green!50!black,
  fonttitle=\bfseries, title={#1}, arc=3pt, breakable
}
\newtcolorbox{pointbox}{
  colback=orange!8, colframe=orange!70!black,
  fonttitle=\bfseries, title={ポイント}, arc=3pt, breakable
}
\newtcolorbox{exbox}[1]{
  colback=gray!6, colframe=gray!50,
  fonttitle=\bfseries, title={例：#1}, arc=3pt, breakable
}
\newtcolorbox{warningbox}{
  colback=red!5, colframe=red!60!black,
  fonttitle=\bfseries, title={注意}, arc=3pt, breakable
}

\title{{\large Hoff 第6章} \\[6pt]
       \textbf{6.3\; Sampling from the Conditional Distributions} \\[4pt]
       {\normalsize （条件付き分布からのサンプリング）ゼミ資料}}
\author{}
\date{}

\begin{document}
\maketitle
\thispagestyle{empty}

% ===========================
\section{この節の狙い}
% ===========================

6.2節の離散近似は，グリッドの置き場所と次元の呪いという代償が大きすぎた．
そこで\textbf{サンプリングの路線に戻る}．ただし作戦を変える．

\medskip
第5章・6.1節の作戦は「\textbf{周辺}分布 $p(\sigma^2 \mid \boldsymbol{y})$ から引いて，
次に条件付き分布 $p(\theta \mid \sigma^2, \boldsymbol{y})$ から引く」であった．
これは半共役事前分布のもとで\textbf{1段目が破綻した}．

\medskip
本節の作戦は\textbf{周辺分布を一切使わない}：
\begin{center}
\textbf{完全条件付き分布だけを使って，同時分布からの標本を作れないか？}
\end{center}

\noindent そのために，まず相方である $p(\sigma^2 \mid \theta, \boldsymbol{y})$ を求める．

% ===========================
\section{$\sigma^2$ の完全条件付き分布}
% ===========================

\paragraph{設定}
いま\textbf{$\theta$ の値を知っている}と仮定する（もちろん本当は知らないが，
一時的にそう仮定する）．このとき精度 $\tilde{\sigma}^2 = 1/\sigma^2$ の
条件付き分布を求めたい．

\medskip
出発点は，条件付き分布が同時分布に比例するという事実である
（$\theta$ と $\boldsymbol{y}$ を固定しているので，
$\tilde{\sigma}^2$ に依存しないものはすべて定数として無視できる）：
\begin{align*}
  p(\tilde{\sigma}^2 \mid \theta, y_1, \ldots, y_n)
  &\propto p(y_1, \ldots, y_n,\, \theta,\, \tilde{\sigma}^2) \\
  &= p(y_1, \ldots, y_n \mid \theta, \tilde{\sigma}^2)\;
     p(\theta \mid \tilde{\sigma}^2)\; p(\tilde{\sigma}^2).
\end{align*}

\begin{pointbox}
\textbf{ここで事前独立が効く}

$\theta$ と $\tilde{\sigma}^2$ が事前独立（半共役事前分布）なら
$p(\theta \mid \tilde{\sigma}^2) = p(\theta)$ となり，
これは $\tilde{\sigma}^2$ に\textbf{依存しない定数}なので落とせる：
\[
  p(\tilde{\sigma}^2 \mid \theta, y_1, \ldots, y_n)
  \;\propto\; p(y_1, \ldots, y_n \mid \theta, \tilde{\sigma}^2)\, p(\tilde{\sigma}^2).
\]

\medskip
第5章の\textbf{共役}事前分布だったら，$p(\theta \mid \sigma^2)$ は
$\mathrm{Normal}(\mu_0, \sigma^2/\kappa_0)$ であって
$\sigma^2$ に依存するので\textbf{落とせない}．
6.1節で「半共役は計算を難しくする」と見たが，
完全条件付き分布に関しては\textbf{むしろ簡単にしている}．
\end{pointbox}

\paragraph{計算}
右辺の2つの因子を，$\tilde{\sigma}^2$ の関数として書き下す．

\medskip
\noindent 尤度（$\mathrm{Normal}(\theta, 1/\tilde{\sigma}^2)$ の積）：
\begin{align*}
  p(y_1, \ldots, y_n \mid \theta, \tilde{\sigma}^2)
  &= \prod_{i=1}^{n} \left(\frac{\tilde{\sigma}^2}{2\pi}\right)^{1/2}
    \exp\!\left\{-\frac{\tilde{\sigma}^2 (y_i - \theta)^2}{2}\right\} \\[4pt]
  &\propto (\tilde{\sigma}^2)^{n/2}
    \exp\!\left\{-\tilde{\sigma}^2 \sum_{i=1}^{n}(y_i - \theta)^2 \big/ 2\right\}.
\end{align*}

\noindent 事前分布（$\mathrm{Gamma}(\nu_0/2,\ \nu_0\sigma_0^2/2)$，第2引数はレート）：
\[
  p(\tilde{\sigma}^2)
  \;\propto\; (\tilde{\sigma}^2)^{\nu_0/2 - 1}
    \exp\!\left\{-\tilde{\sigma}^2\, \nu_0\sigma_0^2 / 2\right\}.
\]

\noindent 掛け合わせると，$\tilde{\sigma}^2$ の\textbf{べき}と\textbf{指数の肩}が
それぞれ足し算になる：
\begin{align*}
  p(\tilde{\sigma}^2 \mid \theta, y_1, \ldots, y_n)
  &\propto \underbrace{(\tilde{\sigma}^2)^{n/2}
     \exp\!\left\{-\tilde{\sigma}^2 \textstyle\sum (y_i - \theta)^2 / 2\right\}}_{\text{尤度}}
   \times \underbrace{(\tilde{\sigma}^2)^{\nu_0/2 - 1}
     \exp\!\left\{-\tilde{\sigma}^2 \nu_0\sigma_0^2 / 2\right\}}_{\text{事前分布}} \\[6pt]
  &= (\tilde{\sigma}^2)^{(\nu_0 + n)/2 - 1}
    \times \exp\!\left\{-\tilde{\sigma}^2
      \times \left[\nu_0\sigma_0^2 + \textstyle\sum (y_i - \theta)^2\right] \big/ 2 \right\}.
\end{align*}

\medskip
これは $x^{a-1}e^{-bx}$ の形，すなわち\textbf{ガンマ密度のカーネル}である．
係数を読み取れば
\[
  a = \frac{\nu_0 + n}{2},
  \qquad
  b = \frac{\nu_0\sigma_0^2 + \sum (y_i - \theta)^2}{2}.
\]

\begin{thmbox}{$\sigma^2$ の完全条件付き分布}
半共役事前分布のもと，
\begin{align*}
  \{\tilde{\sigma}^2 \mid \theta, y_1, \ldots, y_n\}
    &\;\sim\; \mathrm{Gamma}\!\left(\frac{\nu_n}{2},\ \frac{\nu_n \sigma_n^2(\theta)}{2}\right), \\[4pt]
  \text{すなわち}\quad
  \{\sigma^2 \mid \theta, y_1, \ldots, y_n\}
    &\;\sim\; \mathrm{inverse\text{-}gamma}\!\left(\frac{\nu_n}{2},\ \frac{\nu_n \sigma_n^2(\theta)}{2}\right)
\end{align*}
ただし
\[
  \nu_n = \nu_0 + n,
  \qquad
  \sigma_n^2(\theta) = \frac{1}{\nu_n}\left[\nu_0 \sigma_0^2 + n s_n^2(\theta)\right],
  \qquad
  s_n^2(\theta) = \frac{1}{n}\sum_{i=1}^{n} (y_i - \theta)^2.
\]
\end{thmbox}

\paragraph{読み方}
$\nu_n = \nu_0 + n$ は「事前の標本サイズ $\nu_0$ ＋実際の標本サイズ $n$」であり，
$\sigma_n^2(\theta)$ は事前の分散 $\sigma_0^2$ と
データから計算した分散 $s_n^2(\theta)$ を\textbf{自由度で重み付けした加重平均}である：
\[
  \sigma_n^2(\theta) = \frac{\nu_0}{\nu_0 + n}\,\sigma_0^2
    + \frac{n}{\nu_0 + n}\, s_n^2(\theta).
\]
$\theta$ の場合（6.1節）は\textbf{精度}で重み付けした加重平均だったが，
ここでは\textbf{自由度}で重み付けした加重平均になっている．
どちらも「事前情報とデータ情報を足し合わせる」という同じ構図である．

\begin{pointbox}
\textbf{$s_n^2(\theta)$ の分母が $n$ であって $n-1$ でないのはなぜか}

通常の標本分散 $s^2 = \sum(y_i - \bar{y})^2/(n-1)$ で $n-1$ で割るのは，
中心 $\bar{y}$ をデータから推定したぶん自由度を1つ失うからである．
ここでは\textbf{$\theta$ を既知として条件付けている}ので，
中心を推定する必要がなく自由度は失われない．
だから $s_n^2(\theta) = \sum(y_i - \theta)^2/n$ が
（$\theta$ が既知なら）$\sigma^2$ の不偏推定量になる．

\medskip
通常の標本分散との関係は，平方和を分解すると見える：
\[
  \sum_{i=1}^{n}(y_i - \theta)^2
  = \underbrace{\sum_{i=1}^{n}(y_i - \bar{y})^2}_{(n-1)s^2}
  + \underbrace{n(\bar{y} - \theta)^2}_{\theta \text{ が } \bar{y} \text{ からずれた分}}.
\]
つまり条件付けた $\theta$ が $\bar{y}$ から遠いほど $s_n^2(\theta)$ は大きくなり，
$\sigma^2$ の条件付き事後分布は\textbf{大きい値のほうへ引っぱられる}．
データのばらつきを「本当のばらつき」と「中心のずれ」の両方で説明するからである．
\end{pointbox}

% ===========================
\section{役者がそろった}
% ===========================

これで2つの\textbf{完全条件付き分布}（full conditional distribution）
——「その1つ以外のすべてを条件にした分布」——が手に入った：

\begin{center}
\renewcommand{\arraystretch}{1.8}
\begin{tabular}{lll}
\toprule
 & 分布 & サンプリング \\
\midrule
$p(\theta \mid \sigma^2, \boldsymbol{y})$
  & $\mathrm{Normal}(\mu_n, \tau_n^2)$（6.1節）
  & \texttt{rnorm()} で可能 \\
$p(\sigma^2 \mid \theta, \boldsymbol{y})$
  & $\mathrm{inverse\text{-}gamma}(\nu_n/2,\ \nu_n\sigma_n^2(\theta)/2)$（本節）
  & \texttt{rgamma()} で可能 \\
\midrule
$p(\theta, \sigma^2 \mid \boldsymbol{y})$
  & 標準的な分布ではない
  & \textbf{直接は不可能} \\
$p(\sigma^2 \mid \boldsymbol{y})$
  & 標準的な分布ではない（6.1節）
  & \textbf{直接は不可能} \\
\bottomrule
\end{tabular}
\end{center}

\medskip
\noindent 部品は2つとも簡単に引けるのに，欲しい同時分布からは引けない．
そこで Hoff の問いはこうなる：

\begin{center}
\textit{Can we use the full conditional distributions to sample from the joint posterior distribution?}\\
（完全条件付き分布を使って，同時事後分布からサンプリングできないだろうか？）
\end{center}

% ===========================
\section{連鎖のアイデア}
% ===========================

議論の土台になるのは，確率の\textbf{乗法定理}を言い換えただけの次の原理である．

\begin{defbox}{サンプリングの合成（composition）}
$X$ を周辺分布 $p(x)$ からの標本とし，
それを条件として $Y$ を条件付き分布 $p(y \mid X)$ から引く．
このとき，ペア $\{X, Y\}$ は\textbf{同時分布 $p(x,y)$ からの標本}である．
\[
  p(x, y) = p(x) \times p(y \mid x)
\]
そして同時分布からの標本は，片方だけ取り出せば\textbf{その周辺分布からの標本}である．
\end{defbox}

\noindent この2つの事実（合成と周辺化）を交互に適用していく．

\paragraph{いま $\sigma^{2(1)}$ が手元にあると仮定する}
すなわち，周辺事後分布 $p(\sigma^2 \mid \boldsymbol{y})$ からの標本が
1つだけ与えられたとする．すると：

\begin{enumerate}
  \item $\theta^{(1)} \sim p(\theta \mid \sigma^{2(1)}, \boldsymbol{y})$ を引く．
        \\
        合成の原理より，$\{\theta^{(1)}, \sigma^{2(1)}\}$ は
        \textbf{同時分布からの標本}である．

  \item ペアが同時分布からの標本なら，片割れ $\theta^{(1)}$ は
        \textbf{周辺分布 $p(\theta \mid \boldsymbol{y})$ からの標本}である．

  \item そこで $\sigma^{2(2)} \sim p(\sigma^2 \mid \theta^{(1)}, \boldsymbol{y})$ を引く．
        \\
        $\theta^{(1)}$ が周辺分布からの標本で，$\sigma^{2(2)}$ がそれを条件とする
        条件付き分布からの標本なのだから，
        合成の原理より $\{\theta^{(1)}, \sigma^{2(2)}\}$ も
        \textbf{同時分布からの標本}である．

  \item ならば $\sigma^{2(2)}$ は
        \textbf{周辺分布 $p(\sigma^2 \mid \boldsymbol{y})$ からの標本}である．
        これは出発点と同じ状況なので，1に戻って繰り返せる．
\end{enumerate}

\begin{center}
\begin{tikzpicture}[
  node distance=10mm and 24mm,
  box/.style={draw, rounded corners=2pt, minimum height=8mm, inner sep=4pt},
  marg/.style={box, fill=green!10, draw=green!50!black},
  arr/.style={-{Stealth[length=2.5mm]}, thick}
]
\node[marg] (s1) {$\sigma^{2(1)}$};
\node[marg, right=of s1] (t1) {$\theta^{(1)}$};
\node[marg, right=of t1] (s2) {$\sigma^{2(2)}$};
\node[marg, right=of s2] (t2) {$\theta^{(2)}$};
\node[right=8mm of t2] (dots) {$\cdots$};

\draw[arr] (s1) -- node[above, font=\scriptsize] {$p(\theta \mid \sigma^2, \boldsymbol{y})$} (t1);
\draw[arr] (t1) -- node[above, font=\scriptsize] {$p(\sigma^2 \mid \theta, \boldsymbol{y})$} (s2);
\draw[arr] (s2) -- node[above, font=\scriptsize] {$p(\theta \mid \sigma^2, \boldsymbol{y})$} (t2);
\draw[arr] (t2) -- (dots);

\node[below=2mm of s1, font=\scriptsize, green!40!black] {周辺からの標本};
\node[below=2mm of t1, font=\scriptsize, green!40!black] {周辺からの標本};
\node[below=2mm of s2, font=\scriptsize, green!40!black] {周辺からの標本};
\node[below=2mm of t2, font=\scriptsize, green!40!black] {周辺からの標本};
\end{tikzpicture}
\end{center}

\noindent 一度この連鎖に乗ってしまえば，
\textbf{完全条件付き分布を交互に使うだけで}，
周辺分布からの標本・同時分布からの標本がいくらでも作れることになる．
隣り合うどのペアを取っても同時事後分布からの標本である．

% ===========================
\section{残った問題：最初の1個をどうするか}
% ===========================

\begin{warningbox}
\textbf{議論が循環している}

上の議論は「$\sigma^{2(1)}$ が周辺事後分布 $p(\sigma^2 \mid \boldsymbol{y})$
からの標本として与えられたとすれば」から始まっていた．
ところが\textbf{その周辺分布から引けないこと}こそ，
6.1節で判明した出発点の問題である．

\medskip
Hoff の原文も
\textit{``if only we had a $\sigma^{2(1)}$ from which to start''}
（出発点となる $\sigma^{2(1)}$ さえあれば）と，
条件法のまま節を終えている．
\end{warningbox}

\begin{pointbox}
\textbf{この議論が本当に示したこと}

厳密に言えば，本節が示したのは
「同時事後分布からスタートすれば，交互サンプリングを続けても
\textbf{同時事後分布から外れない}」ということである．
これは\textbf{同時事後分布がこの手続きの不動点（定常分布）である}
という主張であって，
\textbf{任意の出発点から同時事後分布に近づく}という主張ではない．

\medskip
後者こそが本当に必要な性質であり，
それが成り立つ（したがって $\sigma^{2(1)}$ を適当に決めてよい）ことは，
6.4節でアルゴリズムを定義したうえで，
6.6節のマルコフ連鎖の理論で保証される．
\end{pointbox}

\paragraph{6.4節での解決}
結論を先取りすると，答えは\textbf{「適当な値から始めてよい」}である．
出発点が正しくなくても，交互サンプリングを繰り返すうちに
分布は同時事後分布に収束する．
この手続きを\textbf{ギブスサンプラー}と呼ぶ．
ただし初期の数回分は出発点の影響を引きずるので，
実務上はそれを捨てる（burn-in）ことになる．

% ===========================
\section{まとめ}
% ===========================

\begin{thmbox}{6.3節のまとめ}
\begin{enumerate}
  \item $\theta$ を既知として条件付けると，
        $\{\sigma^2 \mid \theta, \boldsymbol{y}\} \sim
        \mathrm{inverse\text{-}gamma}(\nu_n/2,\ \nu_n\sigma_n^2(\theta)/2)$
        と\textbf{標準的な分布}になる．
        導出の鍵は，事前独立のおかげで $p(\theta \mid \tilde{\sigma}^2) = p(\theta)$
        が定数として落ち，残りがガンマ密度のカーネルになること．

  \item これで\textbf{2つの完全条件付き分布}
        $p(\theta \mid \sigma^2, \boldsymbol{y})$（正規）と
        $p(\sigma^2 \mid \theta, \boldsymbol{y})$（逆ガンマ）が
        どちらも簡単にサンプリングできるようになった．
        一方，同時分布と周辺分布からは相変わらず直接引けない．

  \item \textbf{合成の原理}（$X \sim p(x)$ かつ $Y \mid X \sim p(y \mid x)$
        ならば $\{X,Y\} \sim p(x,y)$）を交互に適用すると，
        完全条件付き分布だけで同時分布からの標本が次々と作れる．

  \item ただしこの議論は，出発点に
        \textbf{周辺分布からの標本 $\sigma^{2(1)}$ を要求する}という循環を含む．
        示せたのは「同時事後分布が手続きの不動点である」ことまでであり，
        任意の出発点からの収束は6.4節・6.6節に委ねられる．
\end{enumerate}
\end{thmbox}

\end{document}
